\section{\texorpdfstring{\(L^p\) 范数把积分变成距离}{Lp 范数把积分变成距离}}
\label{sec:12-Real-Analysis-Support/07-Lp-Spaces-and-Function-Geometry/01}

向量有长度：\(\norm{x}_2=\sqrt{\sum x_i^2}\)，两个向量的距离就是差的长度。函数呢？两个函数 \(f\) 和 \(g\) 在整个定义域上的"总差异"该怎么度量？

一个直接的想法是把函数值当成分量：在每一点 \(x\) 上，\(f(x)\) 和 \(g(x)\) 都有一个差距。积分把这些逐点差异汇总起来。但汇总的方式不止一种——你可以取绝对值的平均值、均方根、最大值……哪种汇总方式算是"正确"的？

答案取决于你关心什么。

\subsection{不同 \(p\) 强调误差的不同侧面}

设 \((\Omega,\mathcal F,\mu)\) 是测度空间。对 \(1\le p<\infty\)，定义

\[
\norm{f}_p=\left(\int|f|^p\,d\mu\right)^{1/p},
\]

并要求它有限。所有这样的可测函数组成 \(L^p(\mu)\)。当 \(p=\infty\) 时，改用 essential supremum：

\[
\norm{f}_\infty=\inf\{M:|f|\le M\ \text{几乎处处}\},
\]

即"忽略零测集后的最坏值"——函数在绝大多数地方不超过多少。

看一个尖峰例子来感受 \(p\) 的选择有多关键。取 \([0,1]\) 上的函数列

\[
f_n(x)=n\mathbf1_{(0,1/n)}(x).
\]

每个 \(f_n\) 都是一个高 \(n\)、宽 \(1/n\) 的矩形，面积恒为 1。计算三种范数：

\[
\norm{f_n}_1=1,\qquad
\norm{f_n}_2=\sqrt{n},\qquad
\norm{f_n}_\infty=n.
\]

同一个函数，三种范数给出三种完全不同的"大小"。\(L^1\) 只看到总质量——面积恒为 1，所以范数恒为 1。\(L^2\) 部分感受到峰值——平方后高峰被放大，范数按 \(\sqrt{n}\) 增长。\(L^\infty\) 完全被峰值支配——范数就是峰高 \(n\)，随着 \(n\) 增加而飙升。

反过来看一个到处都很平缓的函数：在 \([0,1]\) 上恒等于小常数 \(c\)，三种范数都约为 \(c\)。规律清晰：\(p\) 越大，范数越在意"最高的地方有多高"；\(p\) 越小，越在意"总体有多少"。

选范数不是在选一种计算方法，而是在声明你的评价标准。如果你在意总体偏差，用 \(L^1\) 或 \(L^2\)；如果你在意局部最坏情况——比如工程中的最大应力、控制中的最大超调——用 \(L^\infty\)。同样两个函数之间的"距离"，在不同 \(p\) 下可以相差悬殊。

\subsection{H\"older 与 Minkowski：让积分成为合法的距离}

要成为范数，\(\norm{\cdot}_p\) 必须满足三角不等式。对函数积分而言，这远非显然。Minkowski 不等式说它成立：

\[
\norm{f+g}_p\le\norm{f}_p+\norm{g}_p,
\qquad 1\le p\le\infty.
\]

它保证了 \(L^p\) 是线性赋范空间，而 \((f,g)\mapsto\norm{f-g}_p\) 是合法的度量。两个函数在 \(L^p\) 意义下的距离满足距离应有的所有性质：非负、对称、三角不等式、零距离当且仅当函数"几乎处处相等"（这一点下面再谈）。

另一条不可或缺的工具是 H\"older 不等式：若 \(1/p+1/q=1\)（\(p,q>1\)，或 \((p,q)=(1,\infty)\)），

\[
\int|fg|\,d\mu\le\norm{f}_p\,\norm{g}_q.
\]

\(p=q=2\) 时它退化为你熟悉的 Cauchy--Schwarz 不等式：

\[
\left|\int fg\,d\mu\right|
\le\norm{f}_2\norm{g}_2.
\]

概率论中早已见过它的影子：对零均值随机变量，\(\E[XY]\le\sqrt{\E X^2\,\E Y^2}\)，即相关系数的绝对值不超过 1。那一步证明用的正是 H\"older 不等式的 \(L^2\) 特例。

H\"older 不等式还有一个不那么显眼但很实用的推论：在有限测度空间上，大 \(p\) 控制小 \(p\)。若 \(\mu(\Omega)<\infty\) 且 \(p<r\)，则

\[
\norm{f}_p\le\mu(\Omega)^{1/p-1/r}\norm{f}_r,
\]

所以 \(L^r\subseteq L^p\)——高阶矩有限自动推出低阶矩有限。在概率空间（总测度为 1）上，这意味着只要方差存在，期望就一定存在；只要四阶矩有限，方差就一定有限。但在无限测度空间（比如整个 \(\R\) 上的 Lebesgue 测度）上，这个包含关系不成立——大 \(p\) 控制峰值的行为和小 \(p\) 控制尾部的行为互相独立，两头都不包含对方。

\subsection{几乎处处相等才算同一个点}

到这里有一个需要正视的微妙之处。如果在零测集上修改 \(f\) 的值，积分不变，\(\norm{f}_p\) 也不变。那一个"范数为零"的函数不一定是处处为零的函数——它可以在一个零测集上取非零值，其他地方都是零。

为了让"范数为零 \(\Rightarrow\) 函数为零"成立，\(L^p\) 空间中的元素正式地说是等价类：两个几乎处处相等的函数视为同一个点。你是在和整个等价类打交道，而不是和单个函数打交道。

这个技术处理有一个直击直觉的后果：\(L^p\) 中的"函数在 \(x\) 点的值"没有意义。修改一个点不改变元素，所以你不能合法地问"这个 \(L^2\) 函数在 \(x=0.5\) 处的值是多少"。这和第 13 部分会遇到的再生核 Hilbert 空间（RKHS）形成鲜明对照——在 RKHS 中，点态求值不仅是合法的操作，还是连续泛函。\(L^2\) 和 RKHS 的根本分歧就在这一行定义上。

"\(L^2\) 收敛的函数列可以在每个点上都震荡不收敛"也源于此——\(L^p\) 收敛控制的是积分平均，不控制任何指定点的行为。

\subsection{范数收敛与点态收敛分道扬镳}

函数空间上并行的三种收敛方式各有地盘，互不统属：

\begin{itemize}
\tightlist
\item
  \(\norm{f_n-f}_p\to0\)：\(L^p\) 收敛，控制整体平均误差；
\item
  \(f_n(x)\to f(x)\) 几乎处处：控制绝大多数点上的行为，但不控制平均；
\item
  \(\norm{f_n-f}_\infty\to0\)：一致收敛，同时控制平均和每个点。
\end{itemize}

三个方向可以两两独立。尖峰列 \(f_n=n\mathbf1_{(0,1/n)}\) 几乎处处趋于 0（对任何 \(x>0\)，当 \(n>1/x\) 时函数值为 0），但 \(\norm{f_n}_1=1\) 不降——点态好，范数坏。

反过来，让一个高为 1、宽为 \(1/n\) 的方波在 \([0,1]\) 上不断来回移动。它按 \(L^1\) 范数趋于 0——总质量就是 \(1/n\to0\)；但在每个点上，方波无限次经过，每个点上都有无穷多个 \(n\) 使函数值为 1——范数好，点态坏。函数 \(f_n(x)\) 几乎处处不收敛。

两条道路之间并非毫无桥梁。\(L^p\) 收敛可以推出依测度收敛：对任意 \(\varepsilon>0\)，使 \(|f_n-f|\ge\varepsilon\) 的集合的测度趋于 0。依测度收敛再推出存在一个几乎处处收敛的子列——你不一定能从原序列中得到点态收敛，但一定能挑出一个子列。下一节证明 \(L^p\) 的完备性时，这个子列提取技术正是关键工具——从"按范数 Cauchy"走到"按范数收敛到某个 \(L^p\) 函数"，中间就要靠这个子列来构造极限函数。
